%%%%%%%%%%%%%%%%%%%%%%%%%%%%%%%%%%%%%%%%%%%%%%%%%%%%%%%%%%%%%%%%
\section{Conclusions}
\label{sec:discussion_conclusions}

The objective of this thesis was to develop and evaluate \syssimx{} as a practical framework for heterogeneous hybrid co-simulation.
The work addressed the problem that subsystem models from different tools and modeling paradigms cannot be combined by a single solver when they are available only as executable simulation units.
The thesis therefore focused on the orchestration layer that connects such models into one consistent system-level simulation workflow.

The developed framework provides a shared component abstraction with typed and unit-aware ports, backend wrappers, structural analysis, algebraic-loop handling, master algorithms, hybrid event handling, and runtime model switching.
This answers the research questions on heterogeneous representation, dependency-aware orchestration, hybrid event handling, and switching between alternative subsystem models.
The implementation-level verification in Chapter~\ref{chap:implementation} showed that these mechanisms behave as expected in focused scenarios before they are combined in the case study.

The controlled-pendulum case study demonstrated the combined use of the framework features.
The baseline scenario verified numerical convergence toward a monolithic OpenModelica reference as the macro step size was reduced.
The contact scenario showed that heterogeneous pendulum models, hybrid contact events, PID reset, FEM contact, and runtime model switching can be combined in one closed-loop simulation.
The performance benchmark showed that restricting FEM execution to the contact-relevant part of the trajectory reduced the total median wall time from \(467.6\,\mathrm{s}\) to \(289.6\,\mathrm{s}\) for the selected configuration.

The results support the feasibility of the framework for the tested class of heterogeneous hybrid simulation scenarios.
Section~\ref{sec:discussion_limitations} defines the scope of this claim.

The main contribution is an implementation architecture that combines heterogeneous simulation tools, dependency-aware orchestration, hybrid event handling, and runtime fidelity switching in one controlled co-simulation workflow.
The controlled pendulum served as the evaluation system for this architecture.
